\makeatletter
\let\@oldaddcontentsline\addcontentsline
\renewcommand{\addcontentsline}[3]{}
\section*{附录}
\let\addcontentsline\@oldaddcontentsline
\makeatother

\begin{table}[H]
  \centering
  \caption{支撑材料文件列表}
  \label{tab:appendix}
  \renewcommand{\arraystretch}{1.2}
  \begin{tabularx}{\textwidth}{LX}
    \toprule
    文件 & 说明 \\
    \midrule
    src/io\_data.py & 数据读取与锚点校验 \\
    src/energy\_model.py & 能源结算（无储能解析/LP + 储能 LP-first/MILP） \\
    src/task\_model.py & 任务调度与 GPU 容量约束 \\
    src/forecast.py & 到达需求岭回归预测 \\
    src/q1.py--q4.py & 四个子问题求解入口 \\
    src/scenarios.py & 情景与 $\rho$ 扫描 \\
    main.py & 求解链顶部入口（可复现运行） \\
    figures/*.png & raw/process/result 三类共 17 幅候选图 \\
    \bottomrule
  \end{tabularx}
\end{table}

\subsection*{附录A 源程序代码（核心结算与调度）}

\begin{lstlisting}[language=Python]
# 无储能电力平衡的三段边际成本（简化核心，与正文式 6 对应）
def no_storage_mc(L, ren, sell, buy, export_cap):
    surplus = ren - L
    if surplus > export_cap:        # 弃电区段
        return 0.0
    if surplus > 0:                 # 外送区段
        return sell
    return buy                      # 购电区段

# 逐区域储能 LP-first：目标最小化购电成本 - 外送收益
# 约束：电力平衡 / 购售上限 / SOC 递归 / 端 SOC 守恒
# 完整可运行代码见支撑材料 src/ 目录
\end{lstlisting}
